\scanpage{16}
trouve que dans celle-ci, les $3$-flèches sont triviales (i.e. si source $=$ but, ce sont des identités) - cela un fait qui exprime $H^{0}(B_{G}/X,N)=0$ (i.e. $H^{0}(B_{G},N)\to H^{0}(X,N)$ injectif\ldots). Donc la $3$-catégorie peut être regardée comme une $2$-catégorie - et ``c'est'' celle des Gr-champs sur $X$ épinglés par $G,N$. Si on veut localiser sur $X$, et décrire le $2$-\uline{champ} \inserted{de Picard} sur $X$ des champs de Picard (sur des objets variables de $X$) épinglés par $G,N$, on trouve qu'il est exprime par le complexe
\[
\tau_{\leq 2}\bigl(\mathbb{R}p_{G*}\operatorname{Coker}\bigl(N\to Rq_{G*}\overbrace{C(q_{G}^{*}N)}^{\text{rés. inj.}}\bigr)\bigr)
\]
où $p_{G}:B_{G}\to X$ et $q_{G}:B_{e}\approx X\simeq(B_{G})_{/\unsure{P|p|\rho}}\to B_{G}$.

Toutes ces descriptions étant compatibles avec des variations de $G,N,X$, cela do\unsure{\deleted{nne}|\deleted{nn}e} en principe une description de la $2$-catégorie des Gr-champs, avec $X,G,N$ variables\ldots
% archive sheet number: 15
